\scanpage{55}
\mspage{46}
% archive sheet number: 55

\uline{Homomorphismes \inserted{de \unsure{magasins} et} d'ateliers}

\deleted{Soient}

\uline{Ex 1} Soit $f:{\unsure{L|\mathscr{L}}}'\longrightarrow L$ une application ensembliste, on va en déduire une application canonique
\[
f^*:\mathrm{Figures}(L)\longrightarrow\mathrm{Figures}(L')
\]
Pour toute figure \inserted{\unsure{un|une} $\unsure{\Phi|P}$} dans $L$, on pose
\[
f^*(\Phi)=\{f^{-1}(A)\mid A\in\Phi,\ \underline{f^{-1}(A^\circ)\ne\varnothing}\}\tag*{(1)}
\]
\inserted{i.e. $A^\circ\cap\deleted{\unsure{f_0}}L_0\ne\varnothing$ où $L_0=f(L')$}

On peut considérer $f^*$ comme un composé de \deleted{deux} \inserted{trois} applications ; l'une
% slanted
\margin{\deleted{si $f$ est \unsure{surj.}} \unsure{voir}}
\[
A\longrightarrow\deleted{\unsure{de}}F_{L_0}\bullet A
\]
\margin{$F_{L_0}=\{L_0\}$, figure réduite à la seule strate $\{L_0\}$ si $L_0\ne\varnothing$ i.e. $L'\ne\varnothing$, sinon $F_{L_0}=\varnothing$}
va de
\[
\mathrm{Figures}(\unsure{L'|L})\longrightarrow\mathrm{Figures}(\unsure{L^\circ|L})_{L_0},\tag*{(1 bis)}
\]
le deuxième mbr désigne le sous-ens. du premier, formé des figures à \unsure{support} dans $L_0$. On a d'autre part
\[
\mathrm{Figures}_{L_0}(L)\xrightarrow{\sim}\mathrm{Figures}(L_0)\tag*{(1 ter)}
\]
% slanted
\margin{Le composé de (1 bis) et de (1 ter) est $i^*:\mathrm{Figures}(L)\longrightarrow\mathrm{Figures}(L_0)$ qui \unsure{formellement} \unsure{s'apparente} à un passage au quotient\ldots}
et enfin
\[
f_0^*:\mathrm{Figures}_{\deleted{\unsure{L}}}(L_0)\unsure{\relbar\joinrel\relbar|\longrightarrow}\mathrm{Figures}(L'),
\]
% slanted
\margin{$f_0$ étant l'appl. \unsure{can.} induite $\unsure{L|L'}\longrightarrow L_0$}
où cette fois, vu que $f$ est surjectif, on peut poser simplement
\[
f_0^*(\Phi_0)=\{f^{-1}(A)\mid A\in\Phi_0\},\tag*{(2)}
\]
la condition $f_0^{-1}(A^\circ)\ne\varnothing$ étant automatique
